% Source: https://artofproblemsolving.com/community/c3h1845519_confusing_problem
% Stable identity: aops-topic-1845519-q01
Let the speeds be $v_1,v_2,v_3$ (the first car left first, so it is the slowest) and let $AB=d$.

\textbf{Equal intervals.} The cars take $\dfrac{d}{v_1},\dfrac{d}{v_2},\dfrac{d}{v_3}$ hours to reach $B$ and arrive together, so their departure times differ by the differences of these travel times. Equal intervals mean $\dfrac{d}{v_1}-\dfrac{d}{v_2}=\dfrac{d}{v_2}-\dfrac{d}{v_3}$, i.e.
\[
\dfrac{2}{v_2}=\dfrac{1}{v_1}+\dfrac{1}{v_3}
\]
($v_2$ is the harmonic mean of $v_1$ and $v_3$).

\textbf{The meeting.} All three leave $B$ together. When the first and third cars meet, $80$ miles from $C$, the first car has driven $160$ miles and the third $240+80=320$ miles, in the same time. So $v_3=2v_1$, and then $\dfrac{2}{v_2}=\dfrac{1}{v_1}+\dfrac{1}{2v_1}=\dfrac{3}{2v_1}$, i.e.\ $v_2=\dfrac{4}{3}v_1$.

\textbf{The one-hour gap.} $\dfrac{240}{v_1}-\dfrac{240}{v_2}=1$ gives $\dfrac{240}{v_1}-\dfrac{180}{v_1}=1$, so $v_1=60$. Then $v_3=120$ (and $v_2=80$), and the difference of the speeds of the first and third cars is $120-60=\boxed{\textbf{(A)}\ 60}$ miles per hour.

Check: from $B$, the first car reaches $C$ after $4$ hours and the second after $3$ hours; the third reaches $C$ after $2$ hours, and after $\dfrac{8}{3}$ hours the first car is $160$ miles from $B$ while the third is $240-120\cdot\dfrac{2}{3}=160$ miles from $B$.
